\begin{ex}
    Thống kê điểm kiểm tra giữa kỳ môn Toán của $30$ học sinh lớp 11C được ghi lại ở bảng ghép nhóm sau:
    \begin{center}
        \begin{tabular}{|c|c|c|c|c|c|}
            \hline
            Điểm & $[0;2)$ & $[2;4)$ & $[4;6)$ & $[6;8)$ & $[8;10)$ \\
            \hline
            Số học sinh & $1$& $3$ & $8$ & $11$ & $7$ \\
            \hline
        \end{tabular}
    \end{center}
    Trung vị của mẫu số liệu trên thuộc nhóm nào dưới đây?
    \choice
    {$[2;4)$}
    {$[4;6)$}
    {\True $[6;8)$}
    {$[8;10)$}
    \loigiai{
        Tổng số học sinh là $n=30$. Trung vị $M_e$ là giá trị có thứ tự là $\dfrac{30+1}{2}=15{,}5$.
        \begin{center}
        \begin{tabular}{|c|c|c|c|c|c|}
            \hline
            Điểm & $[0;2)$ & $[2;4)$ & $[4;6)$ & $[6;8)$ & $[8;10)$ \\
            \hline
            Số học sinh & $1$& $3$ & $8$ & $11$ & $7$ \\
            \hline
            Tần số tích lũy  & $1$& $4$ & $12$ & $23$ & $30$ \\
            \hline
        \end{tabular}
    \end{center}
        Dựa vào bảng, ta thấy trung vị thuộc nhóm thứ bốn là nhóm $[6;8)$.
    }
\end{ex}

\begin{ex}
    Trong không gian $Oxyz$, cho mặt phẳng $(P)$ đi qua điểm $M(2;2;1)$ và có một vectơ pháp tuyến $\vec{n}=(5;2;-3)$. Phương trình mặt phẳng $(P)$ là
    \choice
    {$5x+2y-3z-17=0$}
    {$2x+2y+z-11=0$}
    {\True $5x+2y-3z-11=0$}
    {$2x+2y+z-17=0$}
    \loigiai{
        Mặt phẳng $(P)$ đi qua $M(2;2;1)$ và có $\vec{n}=(5;2;-3)$ có phương trình là\\
        $5(x-2)+2(y-2)-3(z-1)=0 \Leftrightarrow 5x+2y-3z-10-4+3=0 \Leftrightarrow 5x+2y-3z-11=0$.
    }
\end{ex}

\begin{ex}
    Cho hàm số $y=f(x)$ liên tục trên $\mathbb{R}$ và có bảng xét dấu của $f'(x)$ như sau:
    \begin{center}
        \begin{tikzpicture}
            \tkzTabInit[nocadre=false,lgt=1.2,espcl=2.5,deltacl=0.6]
            {$x$ /0.6, $f'(x)$ /0.6}
            {$-\infty$, $-2$, $0$, $2$, $+\infty$}
            \tkzTabLine{,+,0,-,0,+,0,+,}
        \end{tikzpicture}
    \end{center}
    Hàm số $y=f(x)$ có bao nhiêu điểm cực trị?
    \choice
    {$4$}
    {\True $2$}
    {$3$}
    {$1$}
    \loigiai{
        Dựa vào bảng xét dấu, $f'(x)$ đổi dấu khi đi qua các điểm $x=-2$ (từ dương sang âm) và $x=0$ (từ âm sang dương).\\
        Tại $x=2$, $f'(x)$ không đổi dấu (vẫn dương).\\
        Vậy hàm số có $2$ điểm cực trị.
    }
\end{ex}
\begin{ex}
    Nguyên hàm của hàm số $f(x)=4^x$ là
    \choice
    {$\dfrac{4^{x+1}}{x+1}+C$}
    {$\dfrac{4^x}{2\ln 2}+C$}
    {\True $\dfrac{4^x}{\ln 4}+C$}
    {$x \cdot 4^{x-1}+C$}
    \loigiai{
        Áp dụng công thức nguyên hàm của hàm số mũ $\displaystyle\int a^x \mathrm{d}x = \dfrac{a^x}{\ln a} + C$.\\
        Ta có $\displaystyle\int 4^x \mathrm{d}x = \dfrac{4^x}{\ln 4} + C$. Vì $\ln 4 = \ln 2^2 = 2\ln 2$ nên kết quả cũng có thể viết là $\dfrac{4^x}{2\ln 2}+C$.
    }
\end{ex}

\begin{ex}
    \immini{Biết hàm số $y=\dfrac{2}{cx+b}$ có đồ thị như hình vẽ bên. Giá trị của $b$ và $c$ là
     \choice
    {$\heva{&b=2\\ &c=-1}$}
    {$\heva{&b=1\\ &c=-1}$}
    {$\heva{&b=2\\ &c=1}$}
    {\True$\heva{&b=-2\\ &c=1}$}}
    {\begin{tikzpicture}[scale=.45,>=stealth, font=\footnotesize, line join=round, line cap=round]
	\def\a{0} \def\b{2} \def\c{1} \def\d{-2} % Hệ số
	\def\xmin{-3} \def\xmax{6}
	\def\ymin{-4} \def\ymax{4}
	%\draw[color=gray!50,dashed] (\xmin,\ymin) grid (\xmax,\ymax);
	\draw[->] (\xmin,0)--(\xmax,0) node [below]{$x$};
	\draw[->] (0,\ymin)--(0,\ymax) node [left]{$y$};
	\clip (\xmin+0.1,\ymin+0.1) rectangle (\xmax-0.1,\ymax-0.1);
	\draw[smooth,samples=300,domain=\xmin:(-\d/\c-0.1)] plot(\x,{(\a*(\x)+\b)/(\c*(\x)+\d)});
	\draw[smooth,samples=300,domain=(-\d/\c+0.1:\xmax)] plot(\x,{(\a*(\x)+\b)/(\c*(\x)+\d)});
	\draw (-\d/\c,\ymin)--(-\d/\c,\ymax);
	\draw (\xmin,\a/\c)--(\xmax,\a/\c);
	\fill (0,0) circle (1pt) node[shift={(135:0.25)}]{$O$};
	\fill (0,-1) circle (1pt) node[shift={(-150:0.3)}]{$-1$};
	\fill (2,0) circle (1pt) node[shift={(-135:0.25)}]{$2$};
\end{tikzpicture}}
    \loigiai{
        Dựa vào đồ thị ta thấy:\\
        1. Đồ thị đi qua điểm $(0;-1)$. Thay vào hàm số ta được: $-1 = \dfrac{2}{c\cdot 0 + b} \Rightarrow b = -2$.\\
        2. Đồ thị có tiệm cận đứng $x=2$. Suy ra nghiệm của mẫu số $cx-2=0$ là $x=2 \Rightarrow c=1$.
    }
\end{ex}
\begin{ex}
    Nghiem của phương trình $3^{2x+1}=27$ là
    \choice
    {$5$}
    {$4$}
    {$2$}
    {\True $1$}
    \loigiai{
        Ta có $3^{2x+1}=27 \Leftrightarrow 3^{2x+1}=3^3 \Leftrightarrow 2x+1=3 \Leftrightarrow 2x=2 \Leftrightarrow x=1$.
    }
\end{ex}

\begin{ex}
    Trong không gian với hệ trục tọa độ $Oxyz$, vectơ nào sau đây là một vectơ pháp tuyến của mặt phẳng $(P)\colon 2x-y+z+3=0$?
    \choice
    {\True $\vec{n}_1=(2;-1;1)$}
    {$\vec{n}_2=(2;1;1)$}
    {$\vec{n}_3=(2;-1;3)$}
    {$\vec{n}_4=(-1;1;3)$}
    \loigiai{
        Dựa vào phương trình mặt phẳng $(P)\colon 2x-y+z+3=0$, ta có một vectơ pháp tuyến là $\vec{n}=(2;-1;1)$.
    }
\end{ex}


\begin{ex}
    \immini{Cho hình chóp $S.ABCD$ có đáy là hình vuông, cạnh bên $SA$ vuông góc với đáy $(ABCD)$. Phát biểu nào sau đây \textbf{sai}?
    \choice
    {\True $CD \perp (SBC)$}
    {$SA \perp (ABC)$}
    {$BC \perp (SAB)$}
    {$BD \perp (SAC)$}}
    {\begin{tikzpicture}[scale=0.6, font=\footnotesize, line join=round, line cap=round, >=stealth]
            \coordinate (A) at (0,2);
            \coordinate (B) at (-1,0);
            \coordinate (C) at (2.5,0);
            \coordinate (D) at ($(A)+(C)-(B)$);
            \coordinate (S) at (0,4);
            \coordinate (O) at ($(A)!0.5!(C)$);

            \draw (S)--(B)--(C)--(D)--(S)--(C);
            \draw[dashed] (S)--(A)--(B) (A)--(D) (A)--(C) (B)--(D);

            \draw[fill=black] (A) circle (1pt) node[left] {$A$};
            \draw[fill=black] (B) circle (1pt) node[left] {$B$};
            \draw[fill=black] (C) circle (1pt) node[below right] {$C$};
            \draw[fill=black] (D) circle (1pt) node[right] {$D$};
            \draw[fill=black] (S) circle (1pt) node[above] {$S$};
            \draw[fill=black] (O) circle (1pt) node[below] {$O$};
        \end{tikzpicture}}
    \loigiai{
        \begin{itemize}
            \item Ta có $SA \perp (ABCD) \Rightarrow SA \perp (ABC)$
            \item Vì $BC \perp AB$ và $BC \perp SA$ nên $BC \perp (SAB)$
            \item Vì $BD \perp AC$ và $BD \perp SA$ nên $BD \perp (SAC)$ 
            \item $CD \perp AD$ và $CD \perp SA$ nên $CD \perp (SAD)$,
            \item  $CD$  không vuông với  $(SBC)$.
        \end{itemize}
    }
\end{ex}
\begin{ex}
    Nếu cấp số nhân $(u_n)$ có số hạng đầu $u_1=3$ và công bội $q=3$ thì số hạng tổng quát của cấp số nhân đó là
    \choice[4]
    {\True $u_n=3^n$}
    {$u_n=3^{n-1}$}
    {$u_n=3^{n+1}$}
    {$u_n=3+3\cdot (n-1)$}
    \loigiai{
        Số hạng tổng quát của cấp số nhân là $u_n = u_1 \cdot q^{n-1} = 3 \cdot 3^{n-1} = 3^n$.
    }
\end{ex}
\begin{ex}
    \immini{Cho hàm số $y=f(x)$ có đồ thị như hình vẽ bên.
       Đường tiệm cận xiên của đồ thị hàm số đã cho là đường thẳng
    \choice
    {$y=x-1$}
    {\True $y=-x-1$}
    {$y=x+1$}
    {$y=-x+1$}}
    {\begin{tikzpicture}[scale=.4,>=stealth, font=\footnotesize, line join=round, line cap=round]
	\def\xmin{-5} \def\xmax{5}
	\def\ymin{-5} \def\ymax{5}
	\def\a{-1} \def\b{-2} \def\c{-2} \def\m{1} \def\n{1} % Hệ số
	%\draw[color=gray!50,dashed] (\xmin,\ymin) grid (\xmax,\ymax);
	\draw[->] (\xmin,0)--(\xmax,0) node [shift=(-110:0.2)]{$x$};
	\draw[->] (0,\ymin)--(0,\ymax) node [shift=(-150:0.2)]{$y$};
	\fill (0,0) circle (1pt) node[shift={(135:0.25)}]{$O$};
	\clip (\xmin+0.1,\ymin+0.1) rectangle (\xmax-0.1,\ymax-0.1);
	\draw[red,smooth,samples=300,domain=-\n/\m+0.01:\xmax] plot(\x,{(\a*(\x)^2+\b*(\x)+\c)/(\m*(\x)+\n)});
	\draw[red,smooth,samples=300,domain=\xmin:-\n/\m-0.01] plot(\x,{(\a*(\x)^2+\b*(\x)+\c)/(\m*(\x)+\n)});
	\draw[] (-\n/\m,\ymin)--(-\n/\m,\ymax);
	\draw[smooth,samples=300,domain=(\xmin:\xmax)] plot(\x,{\a/\m*(\x)+ (\b*\m - \a*\n)/(\m)^2});
	\fill (-1,0) circle (1pt) node[scale=0.8,shift={(-135:0.35)}]{$-1$};
	\fill (0,-1) circle (1pt) node[scale=0.8,shift={(-165:0.25)}]{$-1$};
	\fill (0,-2) circle (1pt) node[scale=0.8,shift={(155:0.25)}]{$-2$};
\end{tikzpicture}}
    \loigiai{
        Dựa vào đồ thị, đường tiệm cận xiên đi qua hai điểm $(0;-1)$ và $(-1;0)$.\\
        Phương trình đường thẳng có dạng $y=ax+b$. Thay tọa độ các điểm ta được $\heva{&b=-1\\ &-a+b=0} \Leftrightarrow \heva{&a=-1\\ &b=-1}$.\\
        Vậy đường tiệm cận xiên là $y=-x-1$.
    }
\end{ex}


\begin{ex}
    Trong không gian $Oxyz$, điểm nào dưới đây thuộc đường thẳng $d \colon \heva{&x=1-t\\ &y=5+t\\ &z=2+3t}$?
    \choice
    {$Q(-1;1;3)$}
    {$P(1;2;5)$}
    {$M(1;1;3)$}
    {\True $N(1;5;2)$}
    \loigiai{
        Thay $t=0$ vào phương trình đường thẳng $d$, ta được $\heva{&x=1\\ &y=5\\ &z=2}$.\\
        Vậy điểm $N(1;5;2)$ thuộc đường thẳng $d$.
    }
\end{ex}

\begin{ex}
    \immini{Cho hình hộp $ABCD.A'B'C'D'$. 
    Khi đó,  $\overrightarrow{A'B}+\overrightarrow{A'D}+\overrightarrow{AA'}$ bằng
       \choice
    {$\overrightarrow{AC'}$}
    {\True $\overrightarrow{A'C}$}
    {$\overrightarrow{AB'}$}
    {$\overrightarrow{AD'}$}}
    {\begin{tikzpicture}[scale=0.5, font=\footnotesize,line join=round, line cap=round, >=stealth]
	\path 
	(0,0) coordinate (A)
	++(-130:2) coordinate (B)
	++(0:5) coordinate (C)
	($(A)+(C)-(B)$) coordinate (D)
	($(A)!1/2!(C)$) coordinate (O)
	;
	\foreach \i in {A,B,C,D}{
		\coordinate (\i') at ($(\i)+(-.6,2.6)$);
	}
	\draw (A')--(B')--(C')--(D')--cycle;
	\draw (B)--(B') (C)--(C') (D)--(D')  (B)--(C)--(D);
	\draw[dashed,thin](B)--(A)--(A') (A)--(D);
		\foreach \i/\g in {A'/90,B'/90,C'/90,D'/90,A/-90,B/-90,C/-90,D/-90}
	\fill[black] (\i) circle(1pt)+(\g:5mm)node[scale=1]{$\i$};
\end{tikzpicture}}
    \loigiai{
       Ta có $\overrightarrow{A'B} + \overrightarrow{AA'} = \overrightarrow{A'B'}$.\\
        Theo quy tắc hình bình hành cho mặt chéo: $\overrightarrow{A'B'} + \overrightarrow{A'D} = \overrightarrow{A'C}$.\\
        Vậy $\overrightarrow{A'B} + \overrightarrow{A'D} + \overrightarrow{AA'} = \overrightarrow{A'C}$.
    }
\end{ex}
\Closesolutionfile{ans}
\Opensolutionfile{ans}[ans/dapan2-DE17]
\phanII
\begin{ex}
	\immini{Một mô hình trò chơi vòng quay ở công viên có chiều cao tối đa $43$ m so với mặt đất, bán kính vòng quay là $20$ m. Hai bạn Hoa và Mai cùng chơi chung lượt quay và ngồi trong hai ca bin $B, C$ mà góc $\widehat{BOC}=90^\circ$ (hình vẽ); $\alpha$ là một góc lượng giác hợp bởi tia đầu $OA$, tia cuối $OB$.}
	{		\begin{tikzpicture}[scale=0.7, font=\footnotesize, line join=round, line cap=round, >=stealth]
			% Hệ trục tọa độ
			\draw[->] (-3,0) -- (4,0) node[below] {$x$};
			\draw[->] (0,-3.2) -- (0,3.5) node[left] {$y$};
			\node[above left] at (0,0) {$O$};
			
			% Vẽ vòng quay
			\draw[thick] (0,0) circle (2.5);
			\draw (0,0) circle (2.4);
			
			% Các nan bánh xe
			\foreach \a in {0,30,...,330} \draw (0,0) -- (\a:2.4);
			
			% Các điểm A, B, C và góc alpha
			\coordinate (A) at (2.4,0);
			\coordinate (B) at (30:2.4);
			\coordinate (C) at (-60:2.4);
			
			\draw[fill=black] (A) circle (1pt) node[shift={(-45:0.25)}]{$A$};
			\draw[fill=black] (B) circle (1pt) node[shift={(45:0.25)}]{$B$};
			\draw[fill=black] (C) circle (1pt) node[shift={(170:0.25)}]{$C$};
			
			% Ký hiệu góc alpha
			\draw (0.5,0) arc (0:30:0.5);
			\node at (15:0.8) {$\alpha$};
			\draw[dashed] (B) -- (2.4,1.2);
			
			% Kích thước bên phải
			\draw[<->] (3.2,0) -- (3.2,2.5) node[midway, right] {$20m$};
			\draw[<->] (3.2,0) -- (3.2,-2.5) node[midway, right] {$20m$};
			\draw[dashed] (0,2.5) -- (3.5,2.5);
			\draw[dashed] (0,-2.5) -- (3.5,-2.5);
			\draw[dashed] (0,-3.2) -- (3.5,-3.2);
			\draw[<->] (3.2,-2.5) -- (3.2,-3.2) node[midway, right] {$3m$};
			\draw[thick] (2.8,-3.2) -- (3.6,-3.2);
	\end{tikzpicture}}
	\choiceTF
	{\True Chiều cao của $B$ so với mặt đất là $h_B = 23 + 20\sin \alpha$ (mét)}
	{\True Khi $\alpha = 45^\circ$ thì chiều cao của $B$ so với mặt đất là $37{,}14$ m (làm tròn kết quả đến hàng phần trăm)}
	{\True Chiều cao của $C$ so với mặt đất là $h_C = 23 - 20\cos \alpha$ (mét)}
	{Khi $B$ ở vị trí có độ cao $33$ m thì $C$ ở độ cao $13$ m so với mặt đất}
	\loigiai{
		Chiều cao tối đa là $43$ m, bán kính $R=20$ m nên tâm $O$ ở độ cao $43-20=23$ m so với mặt đất.\\
		\begin{itemchoice}
			\itemch Tọa độ điểm $B$ là $(20\cos \alpha; 20\sin \alpha)$. Chiều cao của $B$ là $h_B = 23 + y_B = 23 + 20\sin \alpha$.
			\itemch Với $\alpha = 45^\circ$, $h_B = 23 + 20\sin 45^\circ = 23 + 10\sqrt{2} \approx 37{,}14$ m.
			\itemch Vì $\widehat{BOC}=90^\circ$ và $C$ nằm sau $B$ theo chiều quay (âm) nên góc lượng giác của $OC$ là $\alpha - 90^\circ$.\\
			$h_C = 23 + 20\sin(\alpha - 90^\circ) = 23 - 20\cos \alpha$.
			\itemch $h_B = 33 \Rightarrow 23 + 20\sin \alpha = 33 \Rightarrow \sin \alpha = \dfrac{1}{2} \Rightarrow \cos^2 \alpha = \dfrac{3}{4} \Rightarrow \cos \alpha =\pm \dfrac{\sqrt{3}}{2}$ .\\
			Khi đó $h_C = 23 \pm 20 \cdot \dfrac{\sqrt{3}}{2}$ m. \\
			
		\end{itemchoice}
	}
\end{ex}
\begin{ex}
	Cơ quan $Q$ thuê công ty $A$ hoặc công ty $B$ tư vấn dự án với xác suất lần lượt là $0{,}4$ và $0{,}6$. Xác suất dự án phát sinh thêm chi phí khi do công ty $A$ và công ty $B$ tư vấn lần lượt là $0{,}05$ và $0{,}03$.
	\choiceTF
	{\True Xác suất để dự án không phát sinh thêm chi phí là $0{,}962$}
	{\True Biết dự án có phát sinh thêm chi phí, xác suất dự án do công ty $A$ tư vấn là $0{,}53$ (Kết quả làm tròn đến hàng phần trăm)}
	{Nếu dự án không phát sinh thêm chi phí, xác suất dự án do công ty $A$ tư vấn lớn hơn $0{,}4$}
	{\True Nếu cơ quan $Q$ triển khai $3$ dự án độc lập, xác suất có đúng $1$ dự án phát sinh thêm chi phí là $0{,}1$. (Kết quả làm tròn đến hàng phần chục)}
	\loigiai{
		Gọi $E$ là biến cố dự án do công ty $A$ tư vấn, $F$ là biến cố dự án do công ty $B$ tư vấn.\\
		Gọi $X$ là biến cố dự án phát sinh thêm chi phí.\\
		Ta có $P(E)=0{,}4$; $P(F)=0{,}6$; $P(X|E)=0{,}05$; $P(X|F)=0{,}03$.\\
		Xác suất dự án phát sinh thêm chi phí là:\\
		$P(X) = P(E)P(X|E) + P(F)P(X|F) = 0{,}4 \cdot 0{,}05 + 0{,}6 \cdot 0{,}03 = 0{,}038$.
		\begin{itemchoice}
			\itemch Xác suất dự án không phát sinh thêm chi phí là $P(\overline{X}) = 1 - 0{,}038 = 0{,}962$.
			\itemch Xác suất dự án do công ty $A$ tư vấn khi biết dự án phát sinh thêm chi phí là:\\
			$P(E|X) = \dfrac{P(E)P(X|E)}{P(X)} = \dfrac{0{,}4 \cdot 0{,}05}{0{,}038} \approx 0{,}5263 \approx 0{,}53$.
			\itemch Xác suất dự án do công ty $A$ tư vấn khi biết dự án không phát sinh thêm chi phí là:\\
			$P(E|\overline{X}) = \dfrac{P(E)P(\overline{X}|E)}{P(\overline{X})} = \dfrac{0{,}4 \cdot (1-0{,}05)}{0{,}962} = \dfrac{0{,}4 \cdot 0{,}95}{0{,}962} \approx 0{,}395 < 0{,}4$.
			\itemch Gọi $Y$ là số dự án phát sinh thêm chi phí trong $3$ dự án. $Y \sim B(3; 0{,}038)$.\\
			$P(Y=1) = C_3^1 \cdot 0{,}038^1 \cdot (1-0{,}038)^2 = 3 \cdot 0{,}038 \cdot 0{,}962^2 \approx 0{,}1055 \approx 0{,}1$.
		\end{itemchoice}
	}
\end{ex}
\begin{ex}
    Cho một vật đang chuyển động thẳng đều với vận tốc bằng $108$ km/h thì đi vào vùng có lực cản gây ra cho vật một gia tốc $a = -0{,}5v$ (m/s$^2$) với $v$ là vận tốc của vật tính bằng m/s. Trong bài toán ta có $v > 0$ và gốc thời gian tính từ lúc vật bắt đầu đi vào vùng có lực cản.
    \choiceTF
    {\True $\ln|v(t)|$ đạo hàm theo biến thời gian $t$ cho ta biểu thức là $\dfrac{v'(t)}{v(t)}$}
    { Biểu thức vận tốc của vật là: $v(t) = 30 \mathrm{e}^{-0{,}5t}$ (m/s)}
    {Quãng đường vật đi được trong $2$ giây đầu tiên trong vùng lực cản là hơn $38$\,m}
    {\True Với giả thiết vận tốc của vật nhỏ hơn $0{,}1$\,m/s thì vật dừng lại. Tổng quãng đường vật đi được trong vùng lực cản tính theo mét (kết quả làm tròn đến hàng đơn vị) là $60$ m}
    \loigiai{
        Đổi đơn vị vận tốc ban đầu: $v_0 = 108$ km/h $= 30$ m/s.\\
        Gia tốc của vật $a = v'(t) = \dfrac{\mathrm{d}v}{\mathrm{d}t} = -0{,}5v$.\\
        \begin{itemchoice}
            \itemch Ta có $(\ln|v(t)|)' = \dfrac{v'(t)}{v(t)}$.
            \itemch Từ $v'(t) = -0{,}5v(t) \Rightarrow \dfrac{v'(t)}{v(t)} = -0{,}5$.\\
            Lấy nguyên hàm hai vế: $\ln|v(t)| = -0{,}5t + C \Rightarrow v(t) = C' \mathrm{e}^{-0{,}5t}$.\\
            Tại $t=0, v(0) = 30 \Rightarrow C' = 30$. Vậy $v(t) = 30 \mathrm{e}^{-0{,}5t}$ (m/s).
            \itemch Quãng đường vật đi được trong $2$ giây đầu tiên là:
            $$s = \displaystyle\int\limits_0^2 30 \mathrm{e}^{-0{,}5t} \mathrm{\,d}t = \left. -60 \mathrm{e}^{-0{,}5t} \right|_0^2 = -60 \mathrm{e}^{-1} + 60 \approx 37{,}9 \text{ m.}$$
            \itemch Vật dừng lại khi $v(t) = 0{,}1 \Leftrightarrow 30 \mathrm{e}^{-0{,}5t} = 0{,}1 \Leftrightarrow \mathrm{e}^{-0{,}5t} = \dfrac{1}{300} \Leftrightarrow t = 2 \ln 300 \approx 11{,}4$ s.\\
            Tổng quãng đường vật đi được đến khi dừng lại:
            $$S = \displaystyle\int\limits_0^{11{,}4} 30 \mathrm{e}^{-0{,}5t} \mathrm{\,d}t = \left. -60 \mathrm{e}^{-0{,}5t} \right|_0^{11{,}4} = -60 \cdot \dfrac{1}{300} + 60 = 59{,}8 \approx 60 \text{ m.}$$
        \end{itemchoice}
    }
\end{ex}

\begin{ex}
    Trong không gian $Oxyz$, giả sử trục $Ox$ theo hướng Nam, trục $Oy$ theo hướng Đông, trục $Oz$ hướng thẳng đứng lên. Đơn vị trên các trục là km. Một máy bay xuất phát từ $O$, vào thời điểm $9$h$30$ sáng, máy bay đang ở độ cao $9$ km, cách điểm xuất phát theo hướng Nam $150$ km và theo hướng Đông $300$ km. Từ $9$h$30$ trở đi, phi công để chế độ bay tự động, vận tốc riêng của máy bay theo hướng Đông là $750$ km/h, độ cao không đổi. Gió thổi theo hướng Bắc với vận tốc $10$ m/s.
    \begin{center}
        \begin{tikzpicture}[scale=0.7, line join=round, line cap=round, >=stealth]
							\tikzset{maybay/.pic={
								\colorlet{vnablue}{blue!75!cyan}
								\colorlet{bellygray}{gray!30}
								\colorlet{goldline}{yellow!80!orange}
								
								\begin{scope}[xscale=-1, rotate=10, transform shape]
									
									\path[draw=vnablue!50!black, fill=vnablue] (3.0, 2.0) -- (1.5, 3.5) -- (1.0, 3.3) -- (2.5, 1.8) -- cycle;
									
									\path[draw=gray!80!black, fill=gray!50] (11.5, 1.5) -- (7.5, 4.0) -- (7.0, 3.8) -- (10.5, 1.3) -- cycle;
									
									\begin{scope}[shift={(8.8, 0.4)}, rotate=10, scale=0.8]
										\draw[line width=2pt, gray!80!black] (0, 1.0) -- (0, -1.0);
										\draw[line width=2pt, gray!80!black] (-0.8, -1.0) -- (0.8, -1.0);
										\fill[black!90] (-0.8, -1.3) circle (0.25);
										\fill[black!90] (0, -1.3) circle (0.25);
										\fill[black!90] (0.8, -1.3) circle (0.25);
									\end{scope}
									
									\begin{scope}[shift={(8.5, 0.2)}, rotate=10]
										\draw[line width=3pt, gray!80!black] (0, 1.0) -- (0, -1.0);
										\draw[line width=3pt, gray!80!black] (-0.8, -1.0) -- (0.8, -1.0);
										\fill[black!90] (-0.8, -1.3) circle (0.25);
										\fill[black!90] (0, -1.3) circle (0.25);
										\fill[black!90] (0.8, -1.3) circle (0.25);
										\fill[black!90] (-0.6, -1.1) circle (0.25);
										\fill[black!90] (0.2, -1.1) circle (0.25);
										\fill[black!90] (1.0, -1.1) circle (0.25);
									\end{scope}
									
									\begin{scope}[shift={(15.4, 0.5)}]
										\draw[line width=2pt, gray!80!black] (0, 0.5) -- (0, -1.0);
										\fill[black!90] (0, -1.0) circle (0.2);
										\fill[black!90] (0.2, -0.9) circle (0.2);
									\end{scope}
									
									\def\thanmaybay{(15.5,0.29).. controls (15.5,0.29) and (17.72,0.39)..
										(17.85,1.23).. controls (17.91,1.57) and (17.35,1.73)..
										(17.04,1.9).. controls (16.73,2.07) and (16.38,2.34)..
										(16.01,2.47).. controls (15.62,2.62) and (15.19,2.74)..
										(14.77,2.75).. controls (11.33,2.84) and (0.09,2.3)..
										(0.04,1.76).. controls (0.03,1.72) and (0,1.62)..
										(0.03,1.58).. controls (0.4,1.03) and (3.25,0.31)..
										(4.78,0.08).. controls (7.13,-0.26) and (15.5,0.29).. cycle;
									}
									
									\draw[draw=vnablue!50!black, fill=vnablue]\thanmaybay;
									
									\begin{scope}
										\clip \thanmaybay;
										
										\path[draw=vnablue!50!black, fill=black!80] (16.4,2.27).. controls (16.21,2.22) and (15.97,2.19)..
										(15.7,2.16).. controls (15.7,2.07) and (15.71,1.81)..
										(15.75,1.5).. controls (16.25,1.62) and (16.71,1.81)..
										(17.04,1.9);
										\draw[vnablue!50!black] (16.03,2.2) -- (16.24,1.72);
										
										\fill[bellygray] (-0,1.73).. controls (3.69,0.02) and (10.73,0.42)..
										(17.72,1.04) -- (17.47,0.36) -- (4.78,-0.3).. controls (-0.64,0.5) and (-0.05,1.01).. (-0,1.73) -- cycle;
										
										\draw[goldline, line width=1.5pt] (-0,1.73).. controls (3.69,0.02) and (10.73,0.42).. (17.72,1.04);
									\end{scope}
									
									\draw[draw=vnablue!50!black, line width=1pt]\thanmaybay;
									
									\path[draw=vnablue!50!black, fill=vnablue] (1.07,1.87) -- (0.13,4.91).. controls (0.13,4.91) and (0.63,4.99)..
									(0.92,4.97).. controls (1.4,4.95) and (1.98,4.02)..
									(2.72,3.29).. controls (3.23,2.78) and (3.92,2.34)..
									(4.6,2.39);
									
									\begin{scope}[shift={(1.4, 3.6)}, scale=0.45, rotate=-15, xscale=-1]
										\fill[yellow] (0,0) .. controls (-0.3, 0.5) and (-0.4, 1.2) .. (0, 1.6) .. controls (0.4, 1.2) and (0.3, 0.5) .. cycle;
										\fill[yellow] (-0.2, 0.2) .. controls (-0.8, 0.6) and (-1.0, 1.0) .. (-0.8, 1.2) .. controls (-0.4, 1.0) and (-0.2, 0.5) .. cycle;
										\fill[yellow] (0.2, 0.2) .. controls (0.8, 0.6) and (1.0, 1.0) .. (0.8, 1.2) .. controls (0.4, 1.0) and (0.2, 0.5) .. cycle;
										\fill[yellow] (-0.3, 0.1) .. controls (-1.0, 0.3) and (-1.2, 0.6) .. (-1.1, 0.8) .. controls (-0.7, 0.6) and (-0.4, 0.3) .. cycle;
										\fill[yellow] (0.3, 0.1) .. controls (1.0, 0.3) and (1.2, 0.6) .. (1.1, 0.8) .. controls (0.7, 0.6) and (0.4, 0.3) .. cycle;
									\end{scope}
									
									\path[draw=vnablue!50!black, fill=vnablue] (2.81,1.51).. controls (2.78,1.11) and (0.89,1.26)..
									(0.47,1.42).. controls (0.42,1.44) and (0.4,1.57)..
									(0.4,1.57).. controls (1.54,1.69) and (2.84,1.87).. cycle;
									
									\path[draw=gray!80!black, fill=gray!40] (11.61,0.97).. controls (11.64,1.14) and (11.3,1.22)..
									(11.13,1.23).. controls (9.36,1.36) and (7.01,0.95)..
									(5.73,0.71).. controls (5.66,0.7) and (5.61,0.6)..
									(5.66,0.56).. controls (5.75,0.48) and (6.33,0.41)..
									(6.94,0.49).. controls (8.15,0.65) and (9.86,0.81)..
									(10.88,0.69).. controls (11.37,0.63) and (11.59,0.82)..cycle;
									
									\path[draw=gray!80!black, fill=gray!60] (7.49,0.65).. controls (7.73,0.43) and (8.22,0.2)..
									(8.76,0.24).. controls (9.32,0.28) and (10.02,0.54)..(10.27,0.84);
									
									\path[draw=gray!80!black, fill=gray!80] (9.3,0.31).. controls (9.02,0.29) and (8.77,0.26)..
									(8.48,0.17).. controls (8.46,0.03) and (8.41,-0.18)..
									(8.49,-0.27).. controls (8.66,-0.46) and (9.15,-0.42)..(9.37,-0.41);
									
									\path[draw=vnablue!50!black, fill=vnablue] (9.22,0.32).. controls (9.09,0.24) and (9.08,-0.37)..
									(9.28,-0.49).. controls (9.61,-0.7) and (10.57,-0.69)..
									(11,-0.51).. controls (11,-0.51) and (11.25,0.27)..
									(10.99,0.49).. controls (10.84,0.61) and (10.47,0.66)..
									(10.21,0.64).. controls (9.86,0.62) and (9.45,0.48)..cycle;
									
									\begin{scope}[shift={(10.0, 0.05)}, scale=0.15, xscale=-1]
										\fill[yellow] (0,0) .. controls (-0.3, 0.5) and (-0.4, 1.2) .. (0, 1.6) .. controls (0.4, 1.2) and (0.3, 0.5) .. cycle;
										\fill[yellow] (-0.2, 0.2) .. controls (-0.8, 0.6) and (-1.0, 1.0) .. (-0.8, 1.2) .. controls (-0.4, 1.0) and (-0.2, 0.5) .. cycle;
										\fill[yellow] (0.2, 0.2) .. controls (0.8, 0.6) and (1.0, 1.0) .. (0.8, 1.2) .. controls (0.4, 1.0) and (0.2, 0.5) .. cycle;
										\fill[yellow] (-0.3, 0.1) .. controls (-1.0, 0.3) and (-1.2, 0.6) .. (-1.1, 0.8) .. controls (-0.7, 0.6) and (-0.4, 0.3) .. cycle;
										\fill[yellow] (0.3, 0.1) .. controls (1.0, 0.3) and (1.2, 0.6) .. (1.1, 0.8) .. controls (0.7, 0.6) and (0.4, 0.3) .. cycle;
									\end{scope}
									
									\path[shift={(10.9,0.05)}, draw=gray!80!black, fill=gray!30, rotate=5] (0,0) ellipse (0.3 and 0.58);
									\path[shift={(10.9,0.05)}, fill=black!80, rotate=5] (0,0) ellipse (0.15 and 0.4);
									
									\path (4.2, 1.45) coordinate (ta) (13.8, 1.65) coordinate (tb);
									\foreach \i in {0, 0.025, ..., 1}{
										\fill[white, rounded corners=0.5pt] ($(ta)!\i!(tb)$) rectangle +(0.15,0.22);
									}
									
									\path[draw=vnablue!50!black, line width=0.5pt, rounded corners=1pt, fill=vnablue] (14.3,2.05) rectangle (14.7,1.2);
									\path[draw=vnablue!50!black, line width=0.5pt, rounded corners=1pt, fill=vnablue] (8.5, 1.9) rectangle (8.9, 1.1);
									\path[draw=vnablue!50!black, line width=0.5pt, rounded corners=1pt, fill=vnablue] (3.5, 1.8) rectangle (3.9, 1.0);
									
									\node[xscale=-1, text=white, font=\sffamily\bfseries\LARGE] at (12.3, 2.2) {Vietnam Airlines};
									
									\begin{scope}[shift={(15.5, 2.05)}, scale=0.25, xscale=-1]
										\fill[yellow] (0,0) .. controls (-0.3, 0.5) and (-0.4, 1.2) .. (0, 1.6) .. controls (0.4, 1.2) and (0.3, 0.5) .. cycle;
										\fill[yellow] (-0.2, 0.2) .. controls (-0.8, 0.6) and (-1.0, 1.0) .. (-0.8, 1.2) .. controls (-0.4, 1.0) and (-0.2, 0.5) .. cycle;
										\fill[yellow] (0.2, 0.2) .. controls (0.8, 0.6) and (1.0, 1.0) .. (0.8, 1.2) .. controls (0.4, 1.0) and (0.2, 0.5) .. cycle;
										\fill[yellow] (-0.3, 0.1) .. controls (-1.0, 0.3) and (-1.2, 0.6) .. (-1.1, 0.8) .. controls (-0.7, 0.6) and (-0.4, 0.3) .. cycle;
										\fill[yellow] (0.3, 0.1) .. controls (1.0, 0.3) and (1.2, 0.6) .. (1.1, 0.8) .. controls (0.7, 0.6) and (0.4, 0.3) .. cycle;
									\end{scope}		
									\node[xscale=-1, text=white, font=\sffamily\tiny] at (2.5, 2.2) {VN-A149};		
								\end{scope}
						}}
					\draw[->] (0,0) -- (6,0) node[below] {$y$ (Đông)};
					\draw[->] (0,0) -- (0,5) node[right] {$z$};
					\draw[->] (0,0) -- (-2.8,-2.8) node[left] {$x$ (Nam)};
					\fill (0,0) circle (1pt) node[shift={(180:0.25)}]{$O$};
					\coordinate (M) at (-2,-2);
					\coordinate (E) at (4.5,0);
					\coordinate (Z) at (0,4);
					\coordinate (ME) at (2.5,-2);
					\coordinate (P) at (2.5,2);
					\fill[red] (P) circle (2pt);
					\draw[dashed] (M) -- (ME) -- (E);
					\draw[dashed] (ME) -- (P)--(Z);
					\node[left] at (M) {$150$};
					\node[above] at (E) {$300$};
					\node[left] at (Z) {$9$};
					\path (P)pic[xscale=-1,scale=0.2,rotate=0]{maybay};				
				
				\end{tikzpicture}
	
    \end{center}
       \choiceTF
    {\True Thời điểm $9$h$30$ thì máy bay đã đi được đoạn đường $\sqrt{150^2 + 300^2 + 9^2}$ km kể từ khi xuất phát}
    {Biết máy bay xuất phát lúc $8$h$00$, vận tốc trung bình của máy bay từ khi xuất phát đến thời điểm $9$h$30$ là $224,5$ km/h (kết quả làm tròn đến hàng phần chục)}
    {\True Từ $9$h$30$ trở đi, phi công để chế độ bay tự động, với vận tốc theo hướng Đông là $750$ km/h, gió thổi theo hướng Bắc với vận tốc $10$ m/s; sau $1$ giờ tiếp theo, máy bay bay được đoạn đường $\sqrt{750^2 + 36^2}$ km}
    {Tọa độ của máy bay lúc $10$h$30$ là $(186; 1050; 9)$}
    \loigiai{
        \begin{itemchoice}
            \itemch Tại thời điểm $9$h$30$, tọa độ máy bay là $M(150; 300; 9)$. Đoạn đường đi được kể từ gốc $O$ là $OM = \sqrt{150^2 + 300^2 + 9^2}$\,km.
            \itemch Thời gian bay từ $8$h$00$ đến $9$h$30$ là $1{,}5$ giờ.\\
            Vận tốc trung bình $v_{tb} = \frac{OM}{1{,}5} = \frac{\sqrt{112581}}{1{,}5} \approx 223{,}4$\,km/h.
            \itemch Đổi vận tốc gió: $10$ m/s $= 36$ km/h theo hướng Bắc (ngược chiều dương trục $Ox$).\\
            Vận tốc tổng hợp của máy bay là $\vec{v} = (-36; 750; 0)$.\\
            Sau $1$ giờ, quãng đường bay thêm là $|\vec{v}| \cdot 1 = \sqrt{(-36)^2 + 750^2} = \sqrt{750^2 + 36^2}$\,km.
            \itemch Sau $1$ giờ (lúc $10$h$30$), máy bay di chuyển thêm vector $\vec{s} = (1 \cdot (-36); 1 \cdot 750; 1 \cdot 0) = (-36; 750; 0)$.\\
            Tọa độ mới là $(150 - 36; 300 + 750; 9) = (114; 1050; 9)$.
        \end{itemchoice}
    }
\end{ex}
\Closesolutionfile{ans}
\Opensolutionfile{ans}[ans/dapan3-DE17]
\phanIII
\begin{ex}
    Trong một vòng tuyển chọn giọng hát, $40$ thí sinh đều phải thể hiện hai phong cách âm nhạc đó là Dân gian và Nhạc nhẹ. Kết quả đánh giá từ Ban giám khảo cho thấy có $28$ thí sinh đạt yêu cầu phong cách Dân gian, $25$ thí sinh đạt yêu cầu phong cách Nhạc nhẹ và có $5$ thí sinh không đạt yêu cầu ở cả hai phong cách. Ban giám khảo chọn ngẫu nhiên $3$ thí sinh từ danh sách những người đã đạt yêu cầu phong cách Dân gian để tham gia phỏng vấn. Tính xác suất để trong $3$ người được chọn có đúng $2$ người đạt thêm yêu cầu phong cách Nhạc nhẹ (Kết quả làm tròn đến chữ số hàng phần trăm).
    \shortans[oly]{$0{,}46$}
    \loigiai{
        Số thí sinh đạt yêu cầu ít nhất một trong hai phong cách là $40 - 5 = 35$ (thí sinh).\\
        Gọi $A$ là tập hợp thí sinh đạt yêu cầu phong cách Dân gian, $B$ là tập hợp thí sinh đạt yêu cầu phong cách Nhạc nhẹ.\\
        Số thí sinh đạt yêu cầu cả hai phong cách là $n(A \cap B) = n(A) + n(B) - n(A \cup B) = 28 + 25 - 35 = 18$ (thí sinh).\\
        Số thí sinh đạt yêu cầu phong cách Dân gian nhưng không đạt yêu cầu phong cách Nhạc nhẹ là $28 - 18 = 10$ (thí sinh).\\
        Ban giám khảo chọn ngẫu nhiên $3$ thí sinh từ $28$ thí sinh đạt yêu cầu phong cách Dân gian nên số phần tử của không gian mẫu là $n(\Omega) = \mathrm{C}_{28}^3 = 3\,276$.\\
        Gọi $X$ là biến cố trong $3$ người được chọn có đúng $2$ người đạt thêm yêu cầu phong cách Nhạc nhẹ.\\
        Số kết quả thuận lợi cho biến cố $X$ là $n(X) = \mathrm{C}_{18}^2 \cdot \mathrm{C}_{10}^1 = 153 \cdot 10 = 1\,530$.\\
        Xác suất của biến cố $X$ là $P(X) = \dfrac{1\,530}{3\,276} \approx 0{,}46703 \ldots$\\
        Làm tròn đến chữ số hàng phần trăm, ta được $0{,}47$.
    }
\end{ex}
\begin{ex}
    Trong đợt thi thử Tốt nghiệp THPT, một học sinh muốn tạo một mật khẩu đăng nhập hệ thống gồm $5$ kí tự phân biệt, được chọn từ tập hợp $8$ chữ cái $S=\{T, H, I, D, O, M, A, N\}$. Để lấy may mắn, học sinh này muốn mật khẩu của mình bắt buộc phải chứa cụm chữ \lq\lq DO\rq\rq\ (chữ $D$ đứng liền trước chữ $O$). Hỏi học sinh đó có thể tạo được bao nhiêu mật khẩu khác nhau thỏa mãn yêu cầu trên?
    \shortans[oly]{$480$}
    \loigiai{
        Vì cụm chữ \lq\lq DO\rq\rq\ luôn đứng liền nhau nên ta coi cụm này là một phần tử $X$. \\
        Bài toán trở thành chọn thêm $3$ chữ cái từ $6$ chữ cái còn lại trong tập $S \setminus \{D, O\}$ và sắp xếp chúng cùng với phần tử $X$ để tạo thành mật khẩu có $4$ vị trí (do $X$ chiếm $2$ vị trí trong mật khẩu $5$ kí tự ban đầu). \\
        Số cách chọn $3$ chữ cái từ $6$ chữ cái còn lại là $\mathrm{C}_6^3$ cách. \\
        Số cách sắp xếp phần tử $X$ và $3$ chữ cái vừa chọn vào $4$ vị trí là $\mathrm{P}_4 = 4!$ cách. \\
        Vậy số mật khẩu có thể tạo được là $\mathrm{C}_6^3 \cdot 4! = 20 \cdot 24 = 480$ mật khẩu.
    }
\end{ex}
\begin{ex}
    Cho hình chóp $S.ABCD$ có đáy là hình chữ nhật và $SA \perp (ABCD)$. Biết rằng $AB=1$; $AD=2$ và $SA=3$. Gọi $M$ là trung điểm $AB$. Tính khoảng cách giữa hai đường thẳng $SC$ và $DM$ (Kết quả làm tròn đến hàng phần trăm)
    \shortans[oly]{$0{,}87$}
    \loigiai{
        \begin{center}
            \begin{tikzpicture}[scale=1, font=\footnotesize, line join=round, line cap=round, >=stealth]
                \coordinate (A) at (0,0);
                \coordinate (B) at (-1.2,-1.2);
                \coordinate (D) at (3.5,0);
                \coordinate (C) at ($(B)+(D)-(A)$);
                \coordinate (S) at (0,3.5);
                \coordinate (M) at ($(A)!0.5!(B)$);
                \coordinate (N) at ($(D)!0.5!(C)$);
                
                \draw (S)--(B)--(C)--(D)--(S)--(C);
                \draw[dashed] (S)--(A)--(B) (A)--(D) (D)--(M) (S)--(N) (M)--(N);
                
               \foreach \i/\g in {A/180,M/180,B/180,C/-100,D/-10,S/90}
	\fill[black] (\i) circle(1pt)+(\g:3mm)node[scale=1]{$\i$};
            \end{tikzpicture}
        \end{center}
        Chọn hệ trục tọa độ $Oxyz$ sao cho $A(0;0;0)$, $B(1;0;0)$, $D(0;2;0)$, $S(0;0;3)$. \\
        Khi đó $M\left(\dfrac{1}{2};0;0\right)$, $C(1;2;0)$. \\
        Ta có $\overrightarrow{DM} = \left(\dfrac{1}{2};-2;0\right)$, $\overrightarrow{SC} = (1;2;-3)$, $\overrightarrow{SD} = (0;2;-3)$. \\
        $[\overrightarrow{DM}, \overrightarrow{SC}] = (6; \dfrac{3}{2}; 3)$. \\
        $d(DM, SC) = \dfrac{|[\overrightarrow{DM}, \overrightarrow{SC}] \cdot \overrightarrow{SD}|}{|[\overrightarrow{DM}, \overrightarrow{SC}]|} = \dfrac{|0 + 3 - 9|}{\sqrt{6^2 + (1,5)^2 + 3^2}} = \dfrac{6}{\sqrt{47{,}25}} \approx 0{,}87$.\\
       \textbf{ Cách 2:}
       	Gọi $P$ là điểm sao cho $MDCP$ là hình bình hành. Khi đó $CP \parallel DM$. \\
		Suy ra $DM \parallel (SCP) \Rightarrow \mathrm{d}(DM, SC) = \mathrm{d}(DM, (SCP)) = \mathrm{d}(M, (SCP))=\dfrac{1}{2}\mathrm{d}(A, (SCP))$.
        	\begin{center}
			\begin{tikzpicture}[scale=1, font=\footnotesize, line join=round, line cap=round, >=stealth]
				\coordinate (A) at (0,0);
				\coordinate (B) at (-1.2,-1.2);
				\coordinate (D) at (3.5,0);
				\coordinate (C) at ($(B)+(D)-(A)$);
				\coordinate (S) at (0,3.5);
				\coordinate (M) at ($(A)!0.5!(B)$);
				\coordinate (P) at ($(C)+(M)-(D)$); % Kẻ CP // DM
				
				\draw (S)--(B)--(C)--(D)--(S)--(C) (C)--(P)--(S);
				\draw[dashed] (S)--(A)--(B) (A)--(D) (D)--(M) (B)--(P);
				
				\foreach \i/\g in {A/180,M/180,B/180,C/-45,D/0,S/90,P/-90}
				\fill[black] (\i) circle (1pt) +(\g:3mm) node {$\i$};
			\end{tikzpicture}
		\end{center}

    }
\end{ex}
\begin{ex}
    Một cửa hàng dự định nhập về hai loại máy tính $A$ và $B$, giá mỗi chiếc lần lượt là $10$ triệu đồng và $20$ triệu đồng. Số vốn ban đầu không vượt quá $4$ tỉ đồng. Loại $A$ mang lại lợi nhuận $2{,}5$ triệu đồng/máy bán được, loại $B$ mang lại lợi nhuận $4$ triệu đồng/máy bán được. Cửa hàng ước tính tổng nhu cầu hàng tháng không quá $250$ máy. Hỏi lợi nhuận lớn nhất của cửa hàng trong tháng đó là bao nhiêu triệu đồng?
    \shortans[oly]{$850$}
    \loigiai{
        Gọi $x, y$ lần lượt là số lượng máy tính loại $A$ và loại $B$ mà cửa hàng nhập về ($x, y \in \mathbb{N}$).\\
        Theo bài ra, ta có hệ bất phương trình giới hạn nguồn vốn và nhu cầu:
        $$\heva{& 10x + 20y \le 4000 \\ & x + y \le 250 \\ & x \ge 0 \\ & y \ge 0} \Leftrightarrow \heva{& x + 2y \le 400 \quad (1) \\ & x + y \le 250 \quad (2) \\ & x \ge 0 \\ & y \ge 0.}$$
        \begin{center}
            \begin{tikzpicture}[line cap=round,line join=round,font=\footnotesize,>=stealth,scale=1.00]
	% Cần thêm vào preamble: \usetikzlibrary{patterns}
	
	\def\xmin{-1.}
	\def\xmax{6.0}
	\def\ymin{-1.}
	\def\ymax{4.0}
	
	\begin{scope}
		\clip (\xmin,\ymin) rectangle (\xmax,\ymax);
		
		% === TÔ MIỀN KHÔNG NGHIỆM (miền nghiệm để trắng) ===
		% Miền KHÔNG thỏa: x>=0
		\fill[pattern=north east lines, pattern color=blue]
		(0.00,4.00) -- (-1.00,4.00) -- (-1.00,-1.00) -- (0.00,-1.00) -- cycle;
		\fill[opacity=0.3,blue]
		(0.00,4.00) -- (-1.00,4.00) -- (-1.00,-1.00) -- (0.00,-1.00) -- cycle;
		
		% Miền KHÔNG thỏa: y>=0
		\fill[pattern=north west lines, pattern color=red]
		(-1.00,0.00) -- (-1.00,-1.00) -- (6.00,-1.00) -- (6.00,0.00) -- cycle;
		\fill[opacity=0.3,red]
		(-1.00,0.00) -- (-1.00,-1.00) -- (6.00,-1.00) -- (6.00,0.00) -- cycle;
		
		% Miền KHÔNG thỏa: 2x+5y<=10
		\fill[pattern=horizontal lines, pattern color=green!70!black]
		(6.00,-0.40) -- (6.00,4.00) -- (-1.00,4.00) -- (-1.00,2.40) -- cycle;
		\fill[opacity=0.3,green!70!black]
		(6.00,-0.40) -- (6.00,4.00) -- (-1.00,4.00) -- (-1.00,2.40) -- cycle;
		
		% Miền KHÔNG thỏa: x+y<=3
		\fill[pattern=vertical lines, pattern color=orange]
		(4.00,-1.00) -- (6.00,-1.00) -- (6.00,4.00) -- (-1.00,4.00) -- cycle;
		\fill[opacity=0.3,orange]
		(4.00,-1.00) -- (6.00,-1.00) -- (6.00,4.00) -- (-1.00,4.00) -- cycle;
		
		% === VẼ CÁC ĐƯỜNG THẲNG ===
		\draw[blue, thick] (0.00,\ymin) -- (0.00,\ymax);
		
		\draw[red, thick] (\xmin,0.00) -- (\xmax,0.00);
		
		\draw[green!70!black, thick, domain=\xmin:\xmax, samples=2]
		plot (\x, {(10.00 - 2.00*\x) / 5.00});
		
		\draw[orange, thick, domain=\xmin:\xmax, samples=2]
		plot (\x, {(3 - 1*\x) / 1});
		
	\end{scope}
	
	% Lưới tọa độ
	\draw[step=1cm,help lines,dashed,opacity=0.3] (\xmin,\ymin) grid (\xmax,\ymax);
	
	% Trục tọa độ
	\draw[->] (\xmin,0)--(\xmax,0) node[above] {$x$};
	\draw[->] (0,\ymin)--(0,\ymax) node[left] {$y$};
	\fill (0,0) circle (1pt) node[below right] {$O$};
	
	% Các giao điểm
	\fill[black] (0.00,0.00) circle (1.5pt);
	\fill[black] (0.00,2.00) circle (1.5pt)node[left]{$A$};
	\fill[black] (0.00,3.00) circle (1.5pt);
	\fill[black] (5.00,-0.00) circle (1.5pt);
	\fill[black] (3.00,-0.00) circle (1.5pt)node[below]{$C$};
	\fill[black] (1.67,1.33) circle (1.5pt)node[above]{$B$};
	
\end{tikzpicture}
        \end{center}
        Lợi nhuận của cửa hàng được tính bởi công thức: $L(x;y) = 2{,}5x + 4y$ (triệu đồng).\\
        Miền nghiệm của hệ bất phương trình là tứ giác $OABC$ với các đỉnh $O(0;0)$, $A(250;0)$, $B(100;150)$, $C(0;200)$.\\
        Ta tính giá trị của $L(x;y)$ tại các đỉnh:
        \begin{itemize}
            \item $L(0;0) = 0$.
            \item $L(250;0) = 2{,}5 \cdot 250 + 4 \cdot 0 = 625$.
            \item $L(100;150) = 2{,}5 \cdot 100 + 4 \cdot 150 = 850$.
            \item $L(0;200) = 2{,}5 \cdot 0 + 4 \cdot 200 = 800$.
        \end{itemize}
        Vậy lợi nhuận lớn nhất của cửa hàng là $850$ triệu đồng khi nhập $100$ máy loại $A$ và $150$ máy loại $B$.
    }
\end{ex}

\begin{ex}
	\immini{Từ một khúc gỗ hình trụ tròn xoay có bán kính đáy bằng $8$ cm, một người thợ mộc cần cắt dọc theo thân trụ để lấy ra $3$ khối gỗ hình hộp chữ nhật bằng nhau. Xét trên mặt cắt ngang (hình tròn đáy), ba khối gỗ hình hộp chữ nhật này tương ứng với ba hình chữ nhật. Mỗi hình chữ nhật có một cạnh nằm trên một cạnh của tam giác đều $ABC$ nội tiếp đường tròn đáy, hai đỉnh còn lại nằm trên đường tròn (tham khảo hình vẽ). Hỏi diện tích mặt cắt ngang của khối gỗ hình hộp chữ nhật lớn nhất bằng bao nhiêu? (kết quả làm tròn đến hàng phần mười). }
    {\begin{tikzpicture}[scale=.85, font=\footnotesize, line join=round, line cap=round, >=stealth]
			\def\R{3}
			% Vẽ đường tròn đáy
			\draw[fill=orange!40!,thick] (0,0) circle (\R);
			
			% Tọa độ các đỉnh tam giác đều ABC
			\coordinate (A) at (90:\R);
			\coordinate (B) at (210:\R);
			\coordinate (C) at (330:\R);
			
			% Vẽ tam giác ABC
			\draw[dashed] (A)--(B)--(C)--cycle;
			
			% Vẽ 3 hình chữ nhật (minh họa 1 cái và xoay)
			\foreach \rot in {0, 120, 240}{
				\begin{scope}[rotate=\rot]
				\coordinate (M1) at (-0.866*\R*0.5, -0.5*\R);
					\coordinate (M2) at (0.866*\R*0.5, -0.5*\R);
					\coordinate (M3) at (1.2, -2.75); % Điểm trên đường tròn
					\coordinate (M4) at (-1.2, -2.75); % Điểm trên đường tròn
					\draw[thick] (-1.3,-1.5) rectangle (1.3,-2.7);
				\end{scope}
			}
			
			\node[above] at (A) {$A$};
			\node[below left] at (B) {$B$};
			\node[below right] at (C) {$C$};
			\draw[fill=black] (0,0) circle (1pt) node[above] {$O$};
		\end{tikzpicture}}
   

	\shortans[oly]{$23{,}6$}
	\loigiai{
		Gọi đường tròn đáy có tâm $O$, bán kính $R=8$ cm. Tam giác đều $ABC$ nội tiếp có cạnh $a = R\sqrt{3} = 8\sqrt{3}$.\\
		Xét một hình chữ nhật $MNPQ$ có $MN$ nằm trên cạnh $BC$, $P, Q$ nằm trên cung nhỏ $BC$.\\
		Gọi $h$ là khoảng cách từ tâm $O$ đến cạnh $BC$, ta có $h = R\cos 60^\circ = \dfrac{R}{2} = 4$.\\
		Gọi $x$ là khoảng cách từ tâm $O$ đến cạnh $PQ$ ($4 < x < 8$).\\
		Độ dài cạnh $PQ = 2\sqrt{R^2 - x^2} = 2\sqrt{64 - x^2}$.\\
		Độ dài cạnh $MQ = x - h = x - 4$.\\
		Diện tích hình chữ nhật là $S(x) = 2(x - 4)\sqrt{64 - x^2}$.\\
		Xét hàm số $f(x) = (x-4)^2(64-x^2)$ trên $(4;8)$.\\
		$f'(x) = 2(x-4)(64-x^2) + (x-4)^2(-2x) = 2(x-4)(64-x^2 - x^2+4x) = 2(x-4)(-2x^2+4x+64)$.\\
		$f'(x) = 0 \Leftrightarrow -2x^2+4x+64 = 0 \Leftrightarrow x^2-2x-32=0 \Leftrightarrow x = 1 + \sqrt{33} \approx 6{,}74$.\\
		Khi đó $S_{max} = 2(1+\sqrt{33}-4)\sqrt{64-(1+\sqrt{33})^2} \approx 23{,}6$.\\
		Vậy diện tích lớn nhất xấp xỉ $23{,}6$ cm$^2$.
	}
\end{ex}
\begin{ex}%[2H5V3-4] 
	Một quả bóng hình cầu có bán kính $r$ đang được treo trong một góc của tường nhà (hai bờ tường vuông góc với nhau). Một điểm $B$ cố định nằm trên mép hai bờ tường và cách mặt đất $80$ cm, sợi dây treo quả bóng có độ dài $A B=30$ cm và đây cũng là độ dài ngắn nhất nối điểm $B$ với mặt xung quanh của quả bóng. 
\begin{center}
\begin{tikzpicture}[scale=.5,>=stealth]
\coordinate [](L) at (0,10);
	\coordinate [label=below:$O$](O) at (0,0);
	%Vẽ bức tường
\coordinate (C) at ($(L)+(0,1)$);
	\coordinate(D) at ($(C)-(4,2)$);
	\coordinate (F) at ($(D)-(3,-1)$);
	\coordinate (E) at ($(D)-(L)+(O)$);
	\coordinate (G) at ($(F)-(D)+(E)$);
	\coordinate (H) at ($(L)+(6,-2)$);
	\coordinate (I) at ($(H)-(L)+(O)$);
	\coordinate (J) at ($(H)+(4,1)$);
	\coordinate (K) at ($(J)-(H)+(I)$);
	%Tô bức tường
	\fill[color=gray!150!black,smooth] (F) --(D)--(L)--(H)--(J)--(K)--(I)--(O)--(E)--(G) -- cycle;
	
\draw[dashed,->](L)--($(L)+(0,1.)$) node [left]{$z$};
\draw[dashed,->](E)--($(O)! 1.5 ! (E)$) node [below]{$x$};
\draw[dashed,->](I)--($(O)! 1.5 ! (I)$) node [below]{$y$};
\draw[line width=2](F)--(D)--(L)--(H)--(J) (G)--(E)--(O)--(I)--(K) ($(O)+ (0,5.25)$)--(L) (O)--($(O)+ (0,.7)$);
\draw[dashed,line width=2]($(O)+ (0,5.25)$)--($(O)+ (0,.7)$) ;
 \coordinate [label=above left:$A$](A) at (-1,5.5);
 \draw ($(O)- (1,-3)$) circle (2.5);
 \draw [dashed]($(O)- (3.5,-3)$) arc (180:0:{2.5} and {1});
 \draw ($(O)+ (-3.5,3)$) arc (-180:0:{2.5} and {1});
\coordinate [label=right:$B$](B) at ($(L)-(0,1)$);
\draw[line width=1.5](A)--(B);
	\foreach \diem in {A,B,L,O}	\fill (\diem)circle(2pt);	
		\draw (-1.5,7.5) node[]{$30$ cm};
		%Vẽ đế
		\coordinate (M) at ($(O)- (1,2)$);
		\coordinate (P) at ($(M)- (1,0)$);
		\coordinate (Q) at ($(M)+ (1,0)$);
		\coordinate (R) at ($(P)- (0,.5)$);
		\coordinate (S) at ($(Q)- (0,.5)$);
		\foreach \diem in {M}	\fill (\diem)circle(2pt);
		\fill ($(O)- (1,-3)$)circle(1.5pt);
\draw(-1.25,2.75) node [left]{$r$};		
		\fill[pattern=north east lines] (P) --(Q)--(S)--(R) -- cycle;	
		\draw[line width=1.5](P)--(Q) (M)--($(M)+ (0,2.5)$);
		\draw[dashed]($(O)- (1,-3)$)--($(O)- (2.1,-2)$);
		\draw (-1,-1) node[left]{$20$ cm};
		\end{tikzpicture}
		
\end{center}Biết rằng quả bóng tiếp xúc với hai bên bờ tường và điểm thấp nhất của quả bóng cách mặt đất $20$ cm. Hỏi khoảng cách giữa hai điểm tiếp xúc của quả bóng với hai bờ tường là bao nhiêu cm? (\textit{Kết quả làm tròn đến hàng phần chục})
\par\shortans[oly]{26{,}9}
\loigiai{
\immini{Minh họa hình vẽ như hình bên.\\Đề bài yêu cầu tìm đường kính của quả cầu hay tính $2 r$.\\
Ta có thể tính được $I K=r \sqrt{2}$; $D E=H O=20$; $B K=80-20-r=60-r$; $I B=r+30$.\\
Áp dụng định lý Pytogo trong $\triangle B I K$ ta có 
\allowdisplaybreaks
\begin{eqnarray*}
K I^2+B K^2=I B^2 &\Leftrightarrow&(r \sqrt{2})^2+(60-r)^2=(r+30)^2\\&\Leftrightarrow& 2 r^2+3600-120 r+r^2=r^2+60 r+900 \\&\Leftrightarrow& 2 r^2-180 r+2700=0 \\&\Leftrightarrow&\heva{&r=40+15 \sqrt{3} \\& r=40-15 \sqrt{3}.}
\end{eqnarray*}
}{
\begin{tikzpicture}[scale=.7,>=stealth]
\coordinate [label=right:$B$](B) at (0,10);
	\coordinate [label=right:$O$](O) at (0,0);
	\coordinate [label=above:$I$](I) at (-4,5);
	\coordinate [label=left:$E$](E) at (-4,0);
	\coordinate [label=below left:$D$](D) at (-4,2.5);
	\coordinate [label=right:$H$](H) at (0,2.5);
	\coordinate [label=right:$K$](K) at (0,5);
	\coordinate [label=right:$A$](A) at ($(I)! .39 ! (B)$);
	\draw (I) circle (2.5);
\foreach \diem in {K,A,B,O,D,I,E,H}	\fill (\diem)circle(1.5pt);		
\draw[->]($(O)- (0,1)$)--($(B)+ (0,1)$);
\draw[](O)--(E)--(I)--(K) (D)--(H) (I)--(B);
\draw (-1.5,5.5) node[]{$r\sqrt{2}$};
\draw (-4.25,3.75) node[]{$r$};
\draw (-3.25,6.25) node[]{$r$};
		\end{tikzpicture}
}Ta loại trường hợp $r=40+15 \sqrt{3}$ do khoảng cách từ $O$ đến $B$ theo đề bài chỉ bằng $80$ cm.\\
Vậy bán kính của quả cầu $r=40-15 \sqrt{3}$.\\
Hình chiếu của tâm $I(r;r;20+r)$ lên các mặt phẳng tọa độ này chính là các điểm tiếp xúc: $M(0;r;20+r)$; $N(r;0;20+r)$.\\
Khoảng cách giữa hai điểm tiếp xúc là đoạn $MN$:\\
$MN=\sqrt{{{(r-0)}^2}+{{(0-r)}^2}+{{(20+r-(20+r))}^2}}=\sqrt{r^2+r^2}=r\sqrt{2}$.
Thay giá trị $ r=45-15\sqrt{3}$ vào, ta được:\\
$MN=(45-15\sqrt{3})\sqrt{2}=45\sqrt{2}-15\sqrt{6}\approx 26,9\text{ (cm)}$.
}
\end{ex}

